
This work validates that layer-wise weight tokenization preserves sufficient structural signal for architecture family classification, achieving 80\% $\pm$ 4.2\% test accuracy on timm Model Zoo---exceeding the 60\% threshold by 20 percentage points and the 25\% random baseline by 55 points. The tokenization method---flattening weight matrices, zero-padding to fixed length, and applying per-layer z-score normalization---retains architecture-specific patterns (batch normalization statistics, attention weights, kernel structures) necessary for property prediction tasks.

Simple per-layer statistical aggregation (mean, standard deviation, L2 norm) matches transformer-based token embedding performance (both 80\% $\pm$ 4.2\%) despite $10\times$ fewer parameters and no cross-layer modeling. On small datasets (70 training samples), architecture families exhibit strong layer-level separability detectable via local aggregation without requiring global attention mechanisms. Transformer overfitting (100\% train, 73.33\% validation) indicates that model capacity must be calibrated to dataset scale in weight-space learning.

An exploratory test of equivariant graph neural networks (GNNs) for capturing local permutation-symmetric patterns showed directional evidence (20\% symmetry differential) but fell short of the hypothesized threshold (30\%), indicating that either proper equivariant implementations are required or the hypothesis overestimated local sensitivity in weight-space tasks. The complementarity hypothesis---that combining transformer and GNN embeddings improves downstream tasks---remains untested.

\textbf{Future work.} Larger datasets (500--1000 models) may reveal when cross-layer attention outperforms statistical aggregation. Harder downstream tasks (property prediction, backdoor detection, model merging) may require richer representations than family classification. Alternative tokenization strategies (graph-based representations, hierarchical encodings, learned embeddings) should be compared. Cross-domain generalization should be tested on NLP transformers, reinforcement learning policies, and diffusion models. Proper equivariant GNN implementations (E(n)-equivariant layers) should be validated on larger datasets to determine whether local permutation sensitivity in weight-space learning is meaningful.

\textbf{Impact.} This work provides (1) a validated tokenization protocol for weight-space transformers (flatten, pad, per-layer z-score normalization), (2) a strong baseline for future comparisons (per-layer statistical aggregation achieves 80\% on family classification), (3) a decision framework for architecture selection (statistics for small datasets, transformers for large or complex tasks), and (4) an open dataset and evaluation (timm Model Zoo, 70/15/15 splits, reproducible code).

Weight-space learning is a validated paradigm with practical tokenization strategies and rigorous baselines. Future work can build on this foundation to enable model zoo search, neural architecture prediction, backdoor detection, and model synthesis applications.
